\section{Discussion}
\label{sec:discussion}

Looking at the traces themselves rather than at final scores, our diagnostics surface several recurring inefficiencies in current LLM-driven evolutionary code search. A non-trivial share of the search budget is spent re-introducing material the run has already discarded: $\sim\!30\%$ of added lines are byte-identical to previously-deleted ones, and this share grows steadily across the trajectory in $118$ of $121$ runs. Breakthrough events are also not crisp, repeatable artifacts: same-prompt replays almost never reproduce the original program byte-for-byte, yet typically recover a substantial fraction of its score from a \emph{different} program. The trajectory carries the structural gain, while the specific program is one draw from a wider distribution. Lineages back to the seed are short, so most of the per-run budget is spent on branches that do not contribute to the final best. On math benchmarks, a Bayesian-optimization pass over a single intermediate program's exposed knobs often matches or exceeds the run's final-best score, suggesting that on math the parametric refinement evolutionary search performs late in a run is largely substitutable by post-hoc tuning. On ALE, two of four frameworks overfit on at least $30\%$ of their problems. These patterns hold across four frameworks, two languages, and five LLMs.

\paragraph{Implications.}
On math, the BO finding suggests a natural decomposition of the work an evolutionary run does: the structural changes it makes early in a run, and the parametric refinement of those structures, which can often be done post-hoc by hyperparameter tuning of an intermediate program. A practical consequence is that math-benchmark headline scores should be reported alongside the single-program tuning ceiling $f^{\star}_{\mathrm{BO}}(s_0)$ so this decomposition is visible; on ALE, public scores should additionally be paired with a private-test re-score to surface overfitting. For system design, cycling growth and lineage shallowness suggest that interventions which prevent the search from re-doing discarded work (lineage-aware credit assignment, deletion-aware novelty filters, prompting strategies that expose a parent's deletion history) are promising directions to explore, complementary to extending the search budget.

\paragraph{Scope and open questions.}
The trace-level view we develop here opens several natural directions. The single-program tuning ceiling, currently reported on math, can be extended to other domains and evaluator types; replay-based reproducibility can be probed with larger samples and richer perturbations of the local search state; and the dynamics we surface (cycling, lineage shallowness, and the frequency--utility gap across edit categories) can be re-examined under different selection rules, prompting strategies, and underlying model families. Because EvoTrace records full source and replay environments, each of these follow-ups can be posed as a controlled intervention on the same trace, without re-running the original search.
